\subsection{泰勒公式余项的积分公式；高阶原函数}

n 阶分部积分公式（II, p.~60, 公式 (11)）使我们可以把余项 \(\mathbf{r}_{n}(x)\)——即一个在区间 I 上具有规则的 \((n+1)^{th}\) 导数的函数的 n 阶泰勒展开中的余项——用一个积分表出（I, p.~22）；事实上，在 (12) 中把 f 换成 \(f'\)、b 换成 x、\(g(t)\) 换成函数 \((t-x)^{n}/n!\)，即得

\[
\begin{array}{l} \mathbf {f} (x) = \mathbf {f} (a) + \mathbf {f} ^ {\prime} (a) \frac {(x - a)}{1 !} + \mathbf {f} ^ {\prime \prime} (a) \frac {(x - a) ^ {2}}{2 !} + \dots \\ \qquad + \mathbf {f} ^ {(n)} (a) \frac {(x - a) ^ {n}}{n !} + \int_ {a} ^ {x} \mathbf {f} ^ {(n + 1)} (t) \frac {(x - t) ^ {n}}{n !} d t \end{array}\tag{17}
\]

换言之

\[
\mathbf {r} _ {n} (x) = \int_ {a} ^ {x} \mathbf {f} ^ {(n + 1)} (t) \frac {(x - t) ^ {n}}{n !} d t,\tag{18}
\]

这个公式常使我们能得到余项的简单界。

给定区间 I 上的一个规则函数 f，f 的任意一个原函数 g 在 I 上连续，从而自身又有原函数；f 的任意原函数的任意原函数称为 f 的二阶原函数。更一般地，f 的 n - 1 阶原函数的一个原函数称为 f 的 n 阶原函数。对 n 用归纳法立即看出，f 的两个 n 阶原函数之差是一个次数至多等于 n - 1 的多项式（系数在 E 中）。若给定了 f 的一个 n 阶原函数在一点 \(a \in I\) 处的值及其前 n - 1 阶导数在该点的值，它就完全确定了。

特别地，\(\int_{a}^{(n)}\mathbf{f}\) 表示 \(n\) 阶原函数（指 \(\mathbf{f}\) 的）中连同其前 \(n - 1\) 阶导数在点 \(a\) 处一同为零的那一个。把 \(n - 1\) 阶泰勒公式应用于这个原函数就表明：若 \(\mathbf{g} = \int_{a}^{(n)}\mathbf{f}\)，则

\[
\mathbf {g} (x) = \int_ {a} ^ {x} \mathbf {f} (t) \frac {(x - t) ^ {n - 1}}{(n - 1) !} d t\tag{19}
\]

从而把 n 阶原函数的确定化为单独一个积分。

